\documentclass[10pt,fleqn]{article} % Default font size and left-justified equations
\usepackage[%
    pdftitle={Informatique : Tris},
    pdfauthor={Xavier Pessoles}]{hyperref}

%%%%%%%%%%%%%%%%%%%%%%%%%%%%%%%%%%%%%%%%%
% Original author:
% Mathias Legrand (legrand.mathias@gmail.com) with modifications by:
% Vel (vel@latextemplates.com)
% License:
% CC BY-NC-SA 3.0 (http://creativecommons.org\newtheorem{rappelT}{Rappel}[section]/licenses/by-nc-sa/3.0/)
%%%%%%%%%%%%%%%%%%%%%%%%%%%%%%%%%%%%%%%%%

%----------------------------------------------------------------------------------------
%	VARIOUS REQUIRED PACKAGES AND CONFIGURATIONS
%----------------------------------------------------------------------------------------

\usepackage[top=2.5cm,bottom=2cm,left=2cm,right=2cm,headsep=40pt,a4paper]{geometry} % Page margins

\usepackage{graphicx} % Required for including pictures
\graphicspath{{images/}} % Specifies the directory where pictures are stored
\usepackage{float}

\usepackage{lipsum} % Inserts dummy text

\usepackage{tikz} % Required for drawing custom shapes

\usepackage[french]{babel} % English language/hyphenation
\frenchbsetup{StandardLists=true} % Pour éviter la collision babel enumitem pour les listes

\usepackage{enumitem} % Customize lists
\setlist{nolistsep} % Reduce spacing between bullet points and numbered lists

\usepackage{booktabs} % Required for nicer horizontal rules in tables

\usepackage{colortbl} % Couleur dans les tableaux

\usepackage{xcolor} % Required for specifying colors by name
%\definecolor{ocre}{RGB}{243,102,25} % Define the orange color used for highlighting throughout the book
\definecolor{ocre}{RGB}{49,133,156} % Couleur ''bleue''
\definecolor{violetf}{RGB}{112,48,160} % Couleur ''violet''
\usepackage{enumitem}
\usepackage{pifont} % Pour les dinglist
\usepackage{multicol}
\usepackage{array} % Centrage vertical dans les tableaux

%----------------------------------------------------------------------------------------
%	FONTS
%----------------------------------------------------------------------------------------

\usepackage{avant} % Use the Avantgarde font for headings
%\usepackage{times} % Use the Times font for headings
%\usepackage{mathptmx} % Use the Adobe Times Roman as the default text font together with math symbols from the Sym­bol, Chancery and Com­puter Modern fonts
\usepackage[adobe-utopia]{mathdesign}
\usepackage{microtype} % Slightly tweak font spacing for aesthetics
\usepackage[utf8]{inputenc} % Required for including letters with accents
\usepackage[T1]{fontenc} % Use 8-bit encoding that has 256 glyphs

%----------------------------------------------------------------------------------------
%	BIBLIOGRAPHY AND INDEX
%----------------------------------------------------------------------------------------

\usepackage[style=alphabetic,citestyle=numeric,sorting=nyt,sortcites=true,autopunct=true,babel=hyphen,hyperref=true,abbreviate=false,backref=true,backend=biber]{biblatex}
\addbibresource{bibliography.bib} % BibTeX bibliography file
\defbibheading{bibempty}{}

\usepackage{calc} % For simpler calculation - used for spacing the index letter headings correctly
\usepackage{makeidx} % Required to make an index
\makeindex % Tells LaTeX to create the files required for indexing

%----------------------------------------------------------------------------------------
%	MAIN TABLE OF CONTENTS
%----------------------------------------------------------------------------------------

\usepackage{titletoc} % Required for manipulating the table of contents

\setcounter{tocdepth}{2}     % Dans la table des matieres
\setcounter{secnumdepth}{2}

\contentsmargin{0cm} % Removes the default margin

% Part text styling
\titlecontents{part}[0cm]
{\addvspace{20pt}\centering\large\bfseries}
{}
{}
{}

% Chapter text styling
\titlecontents{chapter}[1.25cm] % Indentation
{\addvspace{12pt}\large\sffamily\bfseries} % Spacing and font options for chapters
{\color{ocre!60}\contentslabel[\Large\thecontentslabel]{1.25cm}\color{ocre}} % Chapter number
{\color{ocre}}  
{\color{ocre!60}\normalsize\;\titlerule*[.5pc]{.}\;\thecontentspage} % Page number

% Section text styling
\titlecontents{section}[1.25cm] % Indentation
{\addvspace{3pt}\sffamily\bfseries} % Spacing and font options for sections
{\color{ocre!60}\contentslabel[\thecontentslabel]{1.25cm} \color{ocre}} % Section number
{\color{ocre}}
{\hfill\color{ocre!60}\thecontentspage} % Page number
[]

% Subsection text styling
\titlecontents{subsection}[1.25cm] % Indentation
{\addvspace{1pt}\sffamily\small} % Spacing and font options for subsections
{\contentslabel[\thecontentslabel]{1.25cm}} % Subsection number
{}
{\ \titlerule*[.5pc]{.}\;\thecontentspage} % Page number
[]


% Subsection text styling
\titlecontents{subsubsection}[1.25cm] % Indentation
{\addvspace{1pt}\sffamily\small} % Spacing and font options for subsections
{\contentslabel[\thecontentslabel]{1.25cm}} % Subsection number
{}
{\ \titlerule*[.5pc]{.}\;\thecontentspage} % Page number
[]

% List of figures
\titlecontents{figure}[0em]
{\addvspace{-5pt}\sffamily}
{\thecontentslabel\hspace*{1em}}
{}
{\ \titlerule*[.5pc]{.}\;\thecontentspage}
[]

% List of tables
\titlecontents{table}[0em]
{\addvspace{-5pt}\sffamily}
{\thecontentslabel\hspace*{1em}}
{}
{\ \titlerule*[.5pc]{.}\;\thecontentspage}
[]

%----------------------------------------------------------------------------------------
%	MINI TABLE OF CONTENTS IN PART HEADS
%----------------------------------------------------------------------------------------

% Chapter text styling
\titlecontents{lchapter}[0em] % Indenting
{\addvspace{15pt}\large\sffamily\bfseries} % Spacing and font options for chapters
{\color{ocre}\contentslabel[\Large\thecontentslabel]{1.25cm}\color{ocre}} % Chapter number
{}  
{\color{ocre}\normalsize\sffamily\bfseries\;\titlerule*[.5pc]{.}\;\thecontentspage} % Page number

% Section text styling
\titlecontents{lsection}[0em] % Indenting
{\sffamily\small} % Spacing and font options for sections
{\contentslabel[\thecontentslabel]{1.25cm}} % Section number
{}
{}

% Subsection text styling
\titlecontents{lsubsection}[.5em] % Indentation
{\normalfont\footnotesize\sffamily} % Font settings
{}
{}
{}

%----------------------------------------------------------------------------------------
%	PAGE HEADERS
%----------------------------------------------------------------------------------------

\usepackage{fancyhdr} % Required for header and footer configuration



\pagestyle{fancy}
 \renewcommand{\headrulewidth}{0pt}
 \fancyhead{}
 \fancyhead[L]{%
 \noindent\begin{minipage}[c]{2.6cm}%
 \includegraphics[width=2cm]{png/logo_lycee.png}%
 \end{minipage}}

\fancyhead[C]{\rule{8cm}{.5pt}}

 \fancyhead[R]{%
 \noindent\begin{minipage}[c]{3cm}
 \begin{flushright}
 \footnotesize{\textit{\textsf{\xxtete}}}%
 \end{flushright}
 \end{minipage}
}


\fancyfoot[C]{\rule{12cm}{.5pt}}
\renewcommand{\footrulewidth}{0.2pt}
\fancyfoot[C]{\footnotesize{\bfseries \thepage}}
\fancyfoot[L]{ 
\begin{minipage}[c]{.2\linewidth}
\noindent\footnotesize{{\xxauteur}}
\end{minipage}}


\fancyfoot[R]{\footnotesize{\xxpied}
\ifthenelse{\isodd{\value{page}}}{
\begin{tikzpicture}[overlay]
\node[shape=rectangle, 
      rounded corners = .25 cm,
	  draw= ocre,
	  line width=2pt, 
	  fill = ocre!10,
	  minimum width  = 2.5cm,
	  minimum height = 3cm,] at (\xxposongletx,\xxposonglety) {};
\node at (\xxposonglettext,\xxposonglety) {\rotatebox{90}{\textbf{\large\color{ocre}{\xxonglet}}}};
%{};
\end{tikzpicture}}{}
}
%
%
%
% Removes the header from odd empty pages at the end of chapters
\makeatletter
\renewcommand{\cleardoublepage}{
\clearpage\ifodd\c@page\else
\hbox{}
\vspace*{\fill}
\thispagestyle{empty}
\newpage
\fi}

\fancypagestyle{plain}{%
\fancyhf{} % vide l’en-tête et le pied~de~page.
%\fancyfoot[C]{\bfseries \thepage} % numéro de la page en cours en gras
% et centré en pied~de~page.
\fancyfoot[R]{\footnotesize{\xxpied}}
\fancyfoot[C]{\rule{12cm}{.5pt}}
\renewcommand{\footrulewidth}{0.2pt}
\fancyfoot[C]{\footnotesize{\bfseries \thepage}}
\fancyfoot[L]{ 
\begin{minipage}[c]{.2\linewidth}
\noindent\footnotesize{{\xxauteur}}
\end{minipage}}}



%----------------------------------------------------------------------------------------
%	THEOREM STYLES
%----------------------------------------------------------------------------------------

% Conflit avec la police adobe
%\usepackage{amsmath,amsfonts,amssymb,amsthm} % For math equations, theorems, symbols, etc
\usepackage{amsmath,amsthm}

\newcommand{\intoo}[2]{\mathopen{]}#1\,;#2\mathclose{[}}
\newcommand{\ud}{\mathop{\mathrm{{}d}}\mathopen{}}
\newcommand{\intff}[2]{\mathopen{[}#1\,;#2\mathclose{]}}
%\newtheorem{notation}{Notation}[chapter]
\newtheorem{notation}{Notation}[section]

% Boxed/framed environments
\newtheoremstyle{ocrenumbox}% % Theorem style name
{0pt}% Space above
{0pt}% Space below
{\normalfont}% % Body font
{}% Indent amount
{\small\bf\sffamily\color{ocre}}% % Theorem head font
{\;}% Punctuation after theorem head
{0.25em}% Space after theorem head
{\small\sffamily\color{ocre}\thmname{#1}\nobreakspace\thmnumber%{\@ifnotempty{#1}{}\@upn{#2}}% Theorem text (e.g. Theorem 2.1)
\thmnote{\nobreakspace\the\thm@notefont\sffamily\bfseries\color{black}---\nobreakspace#3.}} % Optional theorem note
\renewcommand{\qedsymbol}{$\blacksquare$}% Optional qed square


% Boite pour les corriges
\newtheoremstyle{correctionbox}% % Theorem style name
{0pt}% Space above
{0pt}% Space below
{\normalfont}% % Body font
{}% Indent amount
{\small\bf\sffamily\color{violet}}% % Theorem head font
{\;}% Punctuation after theorem head
{0.25em}% Space after theorem head
{\small\sffamily\color{ocre}\thmname{#1}\nobreakspace\thmnumber%{\@ifnotempty{#1}{}\@upn{#2}}% Theorem text (e.g. Theorem 2.1)
\thmnote{\nobreakspace\the\thm@notefont\sffamily\bfseries\color{black}---\nobreakspace#3.}} % Optional theorem note
\renewcommand{\qedsymbol}{$\blacksquare$}% Optional qed square



\newtheoremstyle{blacknumex}% Theorem style name
{5pt}% Space above
{5pt}% Space below
{\normalfont}% Body font
{} % Indent amount
{\small\bf\sffamily}% Theorem head font
{\;}% Punctuation after theorem head
{0.25em}% Space after theorem head
{\small\sffamily{\tiny\ensuremath{\blacksquare}}\nobreakspace\thmname{#1}\nobreakspace\thmnumber%{\@ifnotempty{#1}{}\@upn{#2}}% Theorem text (e.g. Theorem 2.1)
\thmnote{\nobreakspace\the\thm@notefont\sffamily\bfseries---\nobreakspace#3.}}% Optional theorem note

\newtheoremstyle{blacknumbox} % Theorem style name
{0pt}% Space above
{0pt}% Space below
{\normalfont}% Body font
{}% Indent amount
{\small\bf\sffamily}% Theorem head font
{\;}% Punctuation after theorem head
{0.25em}% Space after theorem head
{\small\sffamily\thmname{#1}\nobreakspace 
\thmnote{\nobreakspace\the\thm@notefont\sffamily\bfseries---\nobreakspace#3.}}% Optional theorem note

% Non-boxed/non-framed environments
\newtheoremstyle{ocrenum}% % Theorem style name
{5pt}% Space above
{5pt}% Space below
{\normalfont}% % Body font
{}% Indent amount
{\small\bf\sffamily\color{ocre}}% % Theorem head font
{\;}% Punctuation after theorem head
{0.25em}% Space after theorem head
{\small\sffamily\color{ocre}\thmname{#1}\nobreakspace%\thmnumber{\@ifnotempty{#1}{}\@upn{#2}}% Theorem text (e.g. Theorem 2.1)
\thmnote{\nobreakspace\the\thm@notefont\sffamily\bfseries\color{black}---\nobreakspace#3.}} % Optional theorem note
\renewcommand{\qedsymbol}{$\blacksquare$}% Optional qed square
\makeatother

% Environnement pour les titres de parties
\newtheoremstyle{partiebox} 
{0pt}% Space above
{0pt}% Space below
{\normalfont}% Body font
{}% Indent amount
{\small\bf\sffamily}% Theorem head font
{\;}% Punctuation after theorem head
{0.25em}% Space after theorem head




% Defines the theorem text style for each type of theorem to one of the three styles above
\newcounter{dummy} 
\numberwithin{dummy}{section}
\theoremstyle{ocrenumbox}
%\newtheorem{theoremeT}[dummy]{Théorème}
\newtheorem{theoremeT}[dummy]{Théorème}
\newtheorem{resultatT}[dummy]{Résultat}
\newtheorem{savoirT}[dummy]{Savoir}
\newtheorem{methodeT}[dummy]{Méthode}
\newtheorem{objectifT}[dummy]{Objectif}
%\newtheorem{problem}{Problem}[chapter]
\newtheorem{problem}{Problem}[section]
%\newtheorem{exerciseT}{Exercise}[chapter]
\newtheorem{exerciseT}{Exercice}[section]

\theoremstyle{blacknumex}
%\newtheorem{exampleT}{Example}[chapter]
\newtheorem{exempleT}{Exemple}[section]
\newtheorem{termT}{Terminal\\}[section]
\newtheorem{pyT}{Python\\}[section]
\newtheorem{sciT}{Scilab\\}[section]
\newtheorem{pseudoT}{Pseudo Code\\}[section]
\newtheorem{sqlT}{SQL\\}[section]

\theoremstyle{blacknumbox}
%\newtheorem{vocabulary}{Vocabulary}[chapter]
\newtheorem{vocabulary}{Vocabulaire}[section]
%\newtheorem{definitionT}{Definition}[section]
\newtheorem{definitionT}{Définition}[section]
\newtheorem{rappelT}{Rappel}[section]
\newtheorem{demoT}{Démonstration}[section]
\newtheorem{corollaryT}[dummy]{Corollaire}
\newtheorem{hypoT}{Hypothèse(s)}

\theoremstyle{ocrenum}
\newtheorem{proposition}[dummy]{Proposition}

\theoremstyle{partiebox}
\newtheorem{titrepartieT}[]{}
\newtheorem{titrechapitreT}[]{}

\theoremstyle{correctionbox}
\newtheorem{correctionT}[dummy]{\color{violet}{Correction}}

%----------------------------------------------------------------------------------------
%	DEFINITION OF COLORED BOXES
%----------------------------------------------------------------------------------------

\RequirePackage[framemethod=tikz]{mdframed} % Required for creating the theorem, definition, exercise and corollary boxes

% Theorem box
\newmdenv[skipabove=7pt,
skipbelow=7pt,
backgroundcolor=ocre!10,
linecolor=ocre,
innerleftmargin=5pt,
innerrightmargin=5pt,
innertopmargin=5pt,
leftmargin=0cm,
rightmargin=0cm,
innerbottommargin=5pt]{tBox}


% Correction
\newmdenv[skipabove=7pt,
skipbelow=7pt,
backgroundcolor=violet!10,
linecolor=violet,
innerleftmargin=5pt,
innerrightmargin=5pt,
innertopmargin=5pt,
leftmargin=0cm,
rightmargin=0cm,
innerbottommargin=5pt]{coBox}


% Exercise box	  
\newmdenv[skipabove=7pt,
skipbelow=7pt,
rightline=false,
leftline=true,
topline=false,
bottomline=false,
backgroundcolor=ocre!10,
linecolor=ocre,
innerleftmargin=5pt,
innerrightmargin=5pt,
innertopmargin=5pt,
innerbottommargin=5pt,
leftmargin=0cm,
rightmargin=0cm,
linewidth=4pt]{eBox}	

% Definition box
\newmdenv[skipabove=7pt,
skipbelow=7pt,
rightline=false,
leftline=true,
topline=false,
bottomline=false,
backgroundcolor=ocre!10,
linecolor=ocre,
innerleftmargin=5pt,
innerrightmargin=5pt,
innertopmargin=1pt,
leftmargin=0cm,
rightmargin=0cm,
linewidth=4pt,
innerbottommargin=0pt]{dBox}	

% Demonstration box
\newmdenv[skipabove=7pt,
skipbelow=7pt,
rightline=false,
leftline=true,
topline=false,
bottomline=false,
%backgroundcolor=ocre!10,
linecolor=ocre,
innerleftmargin=5pt,
innerrightmargin=5pt,
innertopmargin=0pt,
leftmargin=0cm,
rightmargin=0cm,
linewidth=4pt,
innerbottommargin=0pt]{demoBox}	

% Corollary box
\newmdenv[skipabove=7pt,
skipbelow=7pt,
rightline=false,
leftline=true,
topline=false,
bottomline=false,
linecolor=gray,
backgroundcolor=black!5,
innerleftmargin=5pt,
innerrightmargin=5pt,
innertopmargin=5pt,
leftmargin=0cm,
rightmargin=0cm,
linewidth=4pt,
innerbottommargin=5pt]{cBox}


% Hypothèses
\newmdenv[skipabove=7pt,
skipbelow=7pt,
rightline=false,
leftline=true,
topline=false,
bottomline=false,
linecolor=gray,
backgroundcolor=black!5,
innerleftmargin=5pt,
innerrightmargin=5pt,
innertopmargin=5pt,
leftmargin=0cm,
rightmargin=0cm,
linewidth=4pt,
innerbottommargin=5pt]{hyBox}


% Boite pour le titre de la partie (pBox)
\newmdenv[skipabove=7pt,
skipbelow=7pt,
rightline=true,
leftline=false,
topline=false,
bottomline=false,
linecolor=ocre,
backgroundcolor=none,
innerleftmargin=5pt,
innerrightmargin=5pt,
innertopmargin=5pt,
leftmargin=0cm,
rightmargin=0cm,
linewidth=4pt,
innerbottommargin=5pt]{pBox}

% Boite pour le titre du chapitre (chBox)
\newmdenv[skipabove=7pt,
skipbelow=7pt,
rightline=false,
leftline=true,
topline=false,
bottomline=false,
linecolor=ocre,
%backgroundcolor=black!5,
innerleftmargin=5pt,
innerrightmargin=5pt,
innertopmargin=5pt,
leftmargin=0cm,
rightmargin=0cm,
linewidth=4pt,
innerbottommargin=5pt]{chBox}


% Boite pour les exemples
\newmdenv[skipabove=7pt,
skipbelow=7pt,
rightline=false,
leftline=true,
topline=false,
bottomline=false,
linecolor=gray,
backgroundcolor=white,
innerleftmargin=5pt,
innerrightmargin=5pt,
innertopmargin=5pt,
leftmargin=0cm,
rightmargin=0cm,
linewidth=4pt,
innerbottommargin=5pt]{exBox}

% Boite pour le terminal
\newmdenv[skipabove=7pt,
skipbelow=7pt,
rightline=false,
leftline=true,
topline=false,
bottomline=false,
linecolor=gray,
backgroundcolor=white,
innerleftmargin=5pt,
innerrightmargin=5pt,
innertopmargin=5pt,
leftmargin=0cm,
rightmargin=0cm,
linewidth=4pt,
innerbottommargin=5pt]{termBox}


% Boite pour Python
\newmdenv[skipabove=7pt,
skipbelow=7pt,
rightline=false,
leftline=true,
topline=false,
bottomline=false,
linecolor=gray,
backgroundcolor=white,
innerleftmargin=5pt,
innerrightmargin=5pt,
innertopmargin=5pt,
leftmargin=0cm,
rightmargin=0cm,
linewidth=4pt,
innerbottommargin=5pt]{pyBox}

% Boite pour scilab
\newmdenv[skipabove=7pt,
skipbelow=7pt,
rightline=false,
leftline=true,
topline=false,
bottomline=false,
linecolor=gray,
backgroundcolor=white,
innerleftmargin=5pt,
innerrightmargin=5pt,
innertopmargin=5pt,
leftmargin=0cm,
rightmargin=0cm,
linewidth=4pt,
innerbottommargin=5pt]{sciBox}


% Boite pour pseudo
\newmdenv[skipabove=7pt,
skipbelow=7pt,
rightline=false,
leftline=true,
topline=false,
bottomline=false,
linecolor=gray,
backgroundcolor=white,
innerleftmargin=5pt,
innerrightmargin=5pt,
innertopmargin=5pt,
leftmargin=0cm,
rightmargin=0cm,
linewidth=4pt,
innerbottommargin=5pt]{pseudoBox}

% Boite pour pseudo
\newmdenv[skipabove=7pt,
skipbelow=7pt,
rightline=false,
leftline=true,
topline=false,
bottomline=false,
linecolor=gray,
backgroundcolor=white,
innerleftmargin=5pt,
innerrightmargin=5pt,
innertopmargin=5pt,
leftmargin=0cm,
rightmargin=0cm,
linewidth=4pt,
innerbottommargin=5pt]{sqlBox}


% Creates an environment for each type of theorem and assigns it a theorem text style from the "Theorem Styles" section above and a colored box from above
\newenvironment{theorem}{\begin{tBox}\begin{theoremeT}}{\end{theoremeT}\end{tBox}}
\newenvironment{resultat}{\begin{tBox}\begin{resultatT}}{\end{resultatT}\end{tBox}}
\newenvironment{methode}{\begin{tBox}\begin{methodeT}}{\end{methodeT}\end{tBox}}
\newenvironment{savoir}{\begin{tBox}\begin{savoirT}}{\end{savoirT}\end{tBox}}
\newenvironment{obj}{\begin{tBox}\begin{objectifT}}{\end{objectifT}\end{tBox}}
\newenvironment{corrige}{\begin{coBox}\begin{correctionT}}{\end{correctionT}\end{coBox}}
\newenvironment{exercise}{\begin{eBox}\begin{exerciseT}}{\hfill{\color{ocre}\tiny\ensuremath{\blacksquare}}\end{exerciseT}\end{eBox}}				  
\newenvironment{exercice}{\begin{eBox}\begin{exerciseT}}{\hfill{\color{ocre}\tiny\ensuremath{\blacksquare}}\end{exerciseT}\end{eBox}}				  

\newenvironment{definition}{\begin{dBox}\begin{definitionT}}{\end{definitionT}\end{dBox}}	
\newenvironment{defi}{\begin{dBox}\begin{definitionT}}{\end{definitionT}\end{dBox}}	
\newenvironment{demo}{\begin{demoBox}\begin{demoT}}{\end{demoT}\end{demoBox}}	
\newenvironment{rappel}{\begin{dBox}\begin{rappelT}}{\end{rappelT}\end{dBox}}	
%\newenvironment{exemple}{\begin{exempleT}}{\hfill{\tiny\ensuremath{\blacksquare}}\end{exempleT}}		
\newenvironment{corollary}{\begin{cBox}\begin{corollaryT}}{\end{corollaryT}\end{cBox}}
\newenvironment{hypo}{\begin{hyBox}\begin{hypoT}}{\end{hypoT}\end{hyBox}}	\newenvironment{exemple}{\begin{exBox}\begin{exempleT}}{\hfill{\tiny\ensuremath{\blacksquare}}\end{exempleT}\end{exBox}}	
\newenvironment{titrepartie}{\begin{pBox}\begin{titrepartieT}}{\end{titrepartieT}\end{pBox}}	
\newenvironment{titrechapitre}{\begin{chBox}\begin{titrechapitreT}}{\end{titrechapitreT}\end{chBox}}	

\newenvironment{term}{ \begin{termBox}\begin{termT}}{\end{termT}\end{termBox}}
\newenvironment{py}{ \begin{pyBox}\begin{pyT}}{\end{pyT}\end{pyBox}}
\newenvironment{sci}{ \begin{sciBox}\begin{sciT}}{\end{sciT}\end{sciBox}}
\newenvironment{pseudo}{ \begin{pseudoBox}\begin{pseudoT}}{\end{pseudoT}\end{pseudoBox}}
\newenvironment{envsql}{ \begin{sqlBox}\begin{sqlT}}{\end{sqlT}\end{sqlBox}}


%----------------------------------------------------------------------------------------
%	REMARK ENVIRONMENT
%----------------------------------------------------------------------------------------

\newenvironment{remark}{\par\vspace{10pt}\small % Vertical white space above the remark and smaller font size
\begin{list}{}{
\leftmargin=35pt % Indentation on the left
\rightmargin=25pt}\item\ignorespaces % Indentation on the right
\makebox[-2.5pt]{\begin{tikzpicture}[overlay]
\node[draw=ocre!60,line width=1pt,circle,fill=ocre!25,font=\sffamily\bfseries,inner sep=2pt,outer sep=0pt] at (-15pt,0pt){\textcolor{ocre}{R}};\end{tikzpicture}} % Orange R in a circle
\advance\baselineskip -1pt}{\end{list}\vskip5pt} % Tighter line spacing and white space after remark

\newenvironment{rem}{\par\vspace{10pt}\small % Vertical white space above the remark and smaller font size
\begin{list}{}{
\leftmargin=35pt % Indentation on the left
\rightmargin=25pt}\item\ignorespaces % Indentation on the right
\makebox[-2.5pt]{\begin{tikzpicture}[overlay]
\node[draw=ocre!60,line width=1pt,circle,fill=ocre!25,font=\sffamily\bfseries,inner sep=2pt,outer sep=0pt] at (-15pt,0pt){\textcolor{ocre}{R}};\end{tikzpicture}} % Orange R in a circle
\advance\baselineskip -1pt}{\end{list}\vskip5pt} % Tighter line spacing and white space after remark


\newenvironment{warn}{\par\vspace{10pt}\small % Vertical white space above the remark and smaller font size
\begin{list}{}{
\leftmargin=35pt % Indentation on the left
\rightmargin=25pt}\item\ignorespaces % Indentation on the right
\makebox[-2.5pt]{\begin{tikzpicture}[overlay]
\node[draw=red!60,line width=1pt,circle,fill=red!25,font=\sffamily\bfseries,inner sep=2pt,outer sep=0pt] at (-15pt,0pt){\textcolor{black}{!}};\end{tikzpicture}} % Point d'exclamation dans un cercle
\advance\baselineskip -1pt}{\end{list}\vskip5pt} % Tighter line spacing and white space after remark


%----------------------------------------------------------------------------------------
%	SECTION NUMBERING IN THE MARGIN
%----------------------------------------------------------------------------------------
\setcounter{secnumdepth}{3}
\setcounter{tocdepth}{2}



\makeatletter
\renewcommand{\@seccntformat}[1]{\llap{\textcolor{ocre}{\csname the#1\endcsname}\hspace{1em}}}                    
\renewcommand{\section}{\@startsection{section}{1}{\z@}
{-4ex \@plus -1ex \@minus -.4ex}
{1ex \@plus.2ex }
{\normalfont\large\sffamily\bfseries}}
\renewcommand{\subsection}{\@startsection {subsection}{2}{\z@}
{-3ex \@plus -0.1ex \@minus -.4ex}
{0.5ex \@plus.2ex }
{\normalfont\sffamily\bfseries}}
\renewcommand{\subsubsection}{\@startsection {subsubsection}{3}{\z@}
{-2ex \@plus -0.1ex \@minus -.2ex}
{.2ex \@plus.2ex }
{\normalfont\small\sffamily\bfseries}}                        
\renewcommand\paragraph{\@startsection{paragraph}{4}{\z@}
{-2ex \@plus-.2ex \@minus .2ex}
{.1ex}
{\normalfont\small\sffamily\bfseries}}

%----------------------------------------------------------------------------------------
%	PART HEADINGS
%----------------------------------------------------------------------------------------


%----------------------------------------------------------------------------------------
%	CHAPTER HEADINGS
%----------------------------------------------------------------------------------------

% \newcommand{\thechapterimage}{}%
% \newcommand{\chapterimage}[1]{\renewcommand{\thechapterimage}{#1}}%
% \def\@makechapterhead#1{%
% {\parindent \z@ \raggedright \normalfont
% \ifnum \c@secnumdepth >\m@ne
% \if@mainmatter
% \begin{tikzpicture}[remember picture,overlay]
% \node at (current page.north west)
% {\begin{tikzpicture}[remember picture,overlay]
% \node[anchor=north west,inner sep=0pt] at (0,0) {\includegraphics[width=\paperwidth]{\thechapterimage}};
% \draw[anchor=west] (\Gm@lmargin,-9cm) node [line width=2pt,rounded corners=15pt,draw=ocre,fill=white,fill opacity=0.5,inner sep=15pt]{\strut\makebox[22cm]{}};
% \draw[anchor=west] (\Gm@lmargin+.3cm,-9cm) node {\huge\sffamily\bfseries\color{black}\thechapter. #1\strut};
% \end{tikzpicture}};
% \end{tikzpicture}
% \else
% \begin{tikzpicture}[remember picture,overlay]
% \node at (current page.north west)
% {\begin{tikzpicture}[remember picture,overlay]
% \node[anchor=north west,inner sep=0pt] at (0,0) {\includegraphics[width=\paperwidth]{\thechapterimage}};
% \draw[anchor=west] (\Gm@lmargin,-9cm) node [line width=2pt,rounded corners=15pt,draw=ocre,fill=white,fill opacity=0.5,inner sep=15pt]{\strut\makebox[22cm]{}};
% \draw[anchor=west] (\Gm@lmargin+.3cm,-9cm) node {\huge\sffamily\bfseries\color{black}#1\strut};
% \end{tikzpicture}};
% \end{tikzpicture}
% \fi\fi\par\vspace*{270\p@}}}

%-------------------------------------------

\def\@makeschapterhead#1{%
\begin{tikzpicture}[remember picture,overlay]
\node at (current page.north west)
{\begin{tikzpicture}[remember picture,overlay]
\node[anchor=north west,inner sep=0pt] at (0,0) {\includegraphics[width=\paperwidth]{\thechapterimage}};
\draw[anchor=west] (\Gm@lmargin,-9cm) node [line width=2pt,rounded corners=15pt,draw=ocre,fill=white,fill opacity=0.5,inner sep=15pt]{\strut\makebox[22cm]{}};
\draw[anchor=west] (\Gm@lmargin+.3cm,-9cm) node {\huge\sffamily\bfseries\color{black}#1\strut};
\end{tikzpicture}};
\end{tikzpicture}
\par\vspace*{270\p@}}
\makeatother

%----------------------------------------------------------------------------------------
%	HYPERLINKS IN THE DOCUMENTS
%----------------------------------------------------------------------------------------


\hypersetup{hidelinks,backref=true,pagebackref=true,hyperindex=true,colorlinks=false,breaklinks=true,urlcolor= ocre,bookmarks=true,bookmarksopen=false,pdftitle={Title},pdfauthor={Author}}
\usepackage{bookmark}
\bookmarksetup{
open,
numbered,
addtohook={%
\ifnum\bookmarkget{level}=0 % chapter
\bookmarksetup{bold}%
\fi
\ifnum\bookmarkget{level}=-1 % part
\bookmarksetup{color=ocre,bold}%
\fi
}
}

%----------------------------------------------------------------------------------------
%	
%----------------------------------------------------------------------------------------

\newcommand{\thechapterimage}{}%
\newcommand{\chapterimage}[1]{\renewcommand{\thechapterimage}{#1}}%
\def\@makechapterhead#1{%
{\parindent \z@ \raggedright \normalfont
\begin{tikzpicture}[remember picture,overlay]
\node at (current page.north west)
{\begin{tikzpicture}[remember picture,overlay]
\node[anchor=north west,inner sep=0pt] at (0,0) {\includegraphics[width=\paperwidth]{\thechapterimage}};
%\draw[anchor=west] (\Gm@lmargin,-9cm) node [line width=2pt,rounded corners=15pt,draw=ocre,fill=white,fill opacity=0.5,inner sep=15pt]{\strut\makebox[22cm]{}};
%\draw[anchor=west] (\Gm@lmargin+.3cm,-9cm) node {\huge\sffamily\bfseries\color{black}\thechapter. #1\strut};
\end{tikzpicture}};
\end{tikzpicture}
\par\vspace*{270\p@}
}}


\newcounter{exo}


\makeatletter             
\renewcommand{\subparagraph}{\@startsection{exo}{5}{\z@}%
                                    {-2ex \@plus-.2ex \@minus .2ex}%
                                    {0ex}%               
{\normalfont\bfseries Question \hspace{.7cm} }}
\makeatother
\renewcommand{\thesubparagraph}{\arabic{subparagraph}} 
\makeatletter



%\makeatletter             
%\renewcommand{\subparagraph}{\@startsection{subparagraph}{5}{\z@}%
%                                    {-2ex \@plus-.2ex \@minus .2ex}%
%                                    {0ex}%               
%{\normalfont\bfseries Question \hspace{.7cm} }}
%\makeatother
%\renewcommand{\thesubparagraph}{\arabic{subparagraph}} 
%\makeatletter


%%%% Environnement pour inclure du code
\usepackage{textcomp}
\usepackage[french]{algorithm2e}
\usepackage{listings}
\lstloadlanguages{R}   % pour regler les pb d accent utf8 dans les codes
\lstset{language=R} % pour regler les pb d accent utf8 dans les codes
\renewcommand{\lstlistlistingname}{Listings}
\renewcommand{\lstlistingname}{Listing}

\SetKwBlock{Fonction}{Début Fonction}{Fin Fonction}
\SetKwComment{Comment}{start}{end}

\definecolor{Bleu}{rgb}{0.1,0.1,1.0}
\definecolor{Noir}{rgb}{0,0,0}
\definecolor{Grau}{rgb}{0.5,0.5,0.5}
\definecolor{DunkelGrau}{rgb}{0.15,0.15,0.15}
\definecolor{Hellbraun}{rgb}{0.5,0.25,0.0}
\definecolor{Magenta}{rgb}{1.0,0.0,1.0}
\definecolor{Gris}{gray}{0.5}
\definecolor{Vert}{rgb}{0,0.5,0}
\definecolor{SourceHintergrund}{rgb}{1,1.0,0.95}


\lstnewenvironment{python}[1][]{
\lstset{
%escapeinside={\%*}{*)},
inputencoding=utf8,   % pour regler les pb d accent utf8 dans les codes
extendedchars=true,   % pour regler les pb d accent utf8 dans les codes
language=python,
basicstyle=\ttfamily\footnotesize, 	
stringstyle=\color{red}, 
showstringspaces=false, 
alsoletter={1234567890},
otherkeywords={\ , \}, \{},
keywordstyle=\color{blue},
emph={access,and,break,class,continue,def,del,elif ,else,
except,exec,finally,for,from,global,if,import,in,i s,
lambda,not,or,pass,print,raise,return,try,while},
emphstyle=\color{black}\bfseries,
emph={[2]True, False, None, self},
emphstyle=[2]\color{black},
emph={[3]from, import, as},
emphstyle=[3]\color{blue},
upquote=true,
columns=flexible, % pour empecher d'avoir un espacement mono
morecomment=[s]{"""}{"""},
commentstyle=\color{Hellbraun}\slshape, 
%emph={[4]1, 2, 3, 4, 5, 6, 7, 8, 9, 0},
emphstyle=[4]\color{blue},
literate=*{:}{{\textcolor{blue}:}}{1}
{=}{{\textcolor{blue}=}}{1}
{-}{{\textcolor{blue}-}}{1}
{+}{{\textcolor{blue}+}}{1}
{*}{{\textcolor{blue}*}}{1}
{!}{{\textcolor{blue}!}}{1}
{(}{{\textcolor{blue}(}}{1}
{)}{{\textcolor{blue})}}{1}
{[}{{\textcolor{blue}[}}{1}
{]}{{\textcolor{blue}]}}{1}
{<}{{\textcolor{blue}<}}{1}
{>}{{\textcolor{blue}>}}{1}
{COMPLETER}{{\textcolor{red}COMPLETER}}{1},
literate=%
            {é}{{\'{e}}}1
            {è}{{\`{e}}}1
            {ê}{{\^{e}}}1
            {ë}{{\¨{e}}}1
            {û}{{\^{u}}}1
            {ù}{{\`{u}}}1
            {â}{{\^{a}}}1
            {à}{{\`{a}}}1
            {î}{{\^{i}}}1
            {ç}{{\c{c}}}1
            {Ç}{{\c{C}}}1
            {É}{{\'{E}}}1
            {Ê}{{\^{E}}}1
            {À}{{\`{A}}}1
            {Â}{{\^{A}}}1
            {Î}{{\^{I}}}1, % pour regler les pb d accent utf8 dans les codes
%framexleftmargin=1mm, framextopmargin=1mm, frame=shadowbox, rulesepcolor=\color{blue},#1
%backgroundcolor=\color{SourceHintergrund}, 
%framexleftmargin=1mm, framexrightmargin=1mm, framextopmargin=1mm, frame=single, framerule=1pt, rulecolor=\color{black},#1
}}{}



\lstnewenvironment{scilab}[1][]{
\lstset{
language=scilab,
basicstyle=\sffamily\footnotesize, 	
stringstyle=\color{red}, 
showstringspaces=false, 
alsoletter={1234567890},
otherkeywords={\ , \}, \{},
keywordstyle=\color{blue},
emph={access,and,break,class,continue,def,del,elif ,else,
except,exec,finally,for,from,global,if,import,in,i s,
lambda,not,or,pass,print,raise,return,try,while,Debut},
emphstyle=\color{black}\bfseries,
emph={[2]True, False, None, self},
emphstyle=[2]\color{black},
emph={[3]from, import, as},
emphstyle=[3]\color{blue},
upquote=true,
columns=flexible, % pour empecher d'avoir un espacement mono
morecomment=[s]{"""}{"""},
commentstyle=\color{Hellbraun}\slshape, 
%emph={[4]1, 2, 3, 4, 5, 6, 7, 8, 9, 0},
emphstyle=[4]\color{blue},
literate=*{:}{{\textcolor{blue}:}}{1}
{=}{{\textcolor{blue}=}}{1}
{-}{{\textcolor{blue}-}}{1}
{+}{{\textcolor{blue}+}}{1}
{*}{{\textcolor{blue}*}}{1}
{!}{{\textcolor{blue}!}}{1}
{(}{{\textcolor{blue}(}}{1}
{)}{{\textcolor{blue})}}{1}
{[}{{\textcolor{blue}[}}{1}
{]}{{\textcolor{blue}]}}{1}
{<}{{\textcolor{blue}<}}{1}
{>}{{\textcolor{blue}>}}{1},
%framexleftmargin=1mm, framextopmargin=1mm, frame=shadowbox, rulesepcolor=\color{blue},#1
%backgroundcolor=\color{SourceHintergrund}, 
%framexleftmargin=1mm, framexrightmargin=1mm, framextopmargin=1mm, frame=single, framerule=1pt, rulecolor=\color{black},#1
}}{}


\lstdefinestyle{stylepython}{%
escapeinside={\%*}{*)},
inputencoding=utf8,   % pour regler les pb d accent utf8 dans les codes
extendedchars=true,   % pour regler les pb d accent utf8 dans les codes
language=python,
basicstyle=\sffamily\footnotesize, 	
stringstyle=\color{red}, 
showstringspaces=false, 
alsoletter={1234567890},
otherkeywords={\ , \}, \{},
keywordstyle=\color{blue},
emph={access,and,break,class,continue,def,del,elif ,else,
except,exec,finally,for,from,global,if,import,in,i s,
lambda,not,or,pass,print,raise,return,try,while},
emphstyle=\color{black}\bfseries,
emph={[2]True, False, None, self},
emphstyle=[2]\color{green},
emph={[3]from, import, as},
emphstyle=[3]\color{blue},
upquote=true,
columns=flexible, % pour empecher d'avoir un espacement mono
morecomment=[s]{"""}{"""},
commentstyle=\color{Hellbraun}\slshape, 
%emph={[4]1, 2, 3, 4, 5, 6, 7, 8, 9, 0},
emphstyle=[4]\color{blue},
literate=*{:}{{\textcolor{blue}:}}{1}
{=}{{\textcolor{blue}=}}{1}
{-}{{\textcolor{blue}-}}{1}
{+}{{\textcolor{blue}+}}{1}
{*}{{\textcolor{blue}*}}{1}
{!}{{\textcolor{blue}!}}{1}
{(}{{\textcolor{blue}(}}{1}
{)}{{\textcolor{blue})}}{1}
{[}{{\textcolor{blue}[}}{1}
{]}{{\textcolor{blue}]}}{1}
{<}{{\textcolor{blue}<}}{1}
{>}{{\textcolor{blue}>}}{1}
{COMPLETER}{{\textcolor{red}COMPLETER}}{1},
literate=%
            {é}{{\'{e}}}1
            {è}{{\`{e}}}1
            {ê}{{\^{e}}}1
            {ë}{{\¨{e}}}1
            {û}{{\^{u}}}1
            {ù}{{\`{u}}}1
            {â}{{\^{a}}}1
            {à}{{\`{a}}}1
            {î}{{\^{i}}}1
            {ç}{{\c{c}}}1
            {Ç}{{\c{C}}}1
            {É}{{\'{E}}}1
            {Ê}{{\^{E}}}1
            {À}{{\`{A}}}1
            {Â}{{\^{A}}}1
            {Î}{{\^{I}}}1,
%numbers=left,                    % where to put the line-numbers; possible values are (none, left, right)
%numbersep=5pt,                   % how far the line-numbers are from the code
%numberstyle=\tiny\color{mygray}, % the style that is used for the line-numbers
}



\lstnewenvironment{termi}[1][]{
\lstset{
language=scilab,
basicstyle=\sffamily\footnotesize, 	
stringstyle=\color{red}, 
showstringspaces=false, 
alsoletter={1234567890},
otherkeywords={\ , \}, \{},
keywordstyle=\color{blue},
emph={access,and,break,class,continue,def,del,elif ,else,
except,exec,finally,for,from,global,if,import,in,i s,
lambda,not,or,pass,print,raise,return,try,while,Debut},
emphstyle=\color{black}\bfseries,
emph={[2]True, False, None, self},
emphstyle=[2]\color{green},
emph={[3]from, import, as},
emphstyle=[3]\color{blue},
upquote=true,
columns=flexible, % pour empecher d'avoir un espacement mono
morecomment=[s]{"""}{"""},
commentstyle=\color{Hellbraun}\slshape, 
%emph={[4]1, 2, 3, 4, 5, 6, 7, 8, 9, 0},
emphstyle=[4]\color{blue},
literate=*{:}{{\textcolor{blue}:}}{1}
{=}{{\textcolor{blue}=}}{1}
{-}{{\textcolor{blue}-}}{1}
{+}{{\textcolor{blue}+}}{1}
{*}{{\textcolor{blue}*}}{1}
{!}{{\textcolor{blue}!}}{1}
{(}{{\textcolor{blue}(}}{1}
{)}{{\textcolor{blue})}}{1}
{[}{{\textcolor{blue}[}}{1}
{]}{{\textcolor{blue}]}}{1}
{<}{{\textcolor{blue}<}}{1}
{>}{{\textcolor{blue}>}}{1},
%framexleftmargin=1mm, framextopmargin=1mm, frame=shadowbox, rulesepcolor=\color{blue},#1
%backgroundcolor=\color{SourceHintergrund}, 
%framexleftmargin=1mm, framexrightmargin=1mm, framextopmargin=1mm, frame=single, framerule=1pt, rulecolor=\color{black},#1
}}{}


\lstnewenvironment{sql}[1][]{
\lstset{
%escapeinside={\%*}{*)},
%inputencoding=utf8,   % pour regler les pb d accent utf8 dans les codes
%extendedchars=true,   % pour regler les pb d accent utf8 dans les codes
language=sql,
basicstyle=\sffamily\footnotesize, 	
stringstyle=\color{red}, 
showstringspaces=false, 
alsoletter={1234567890},
otherkeywords={\ , \}, \{},
keywordstyle=\color{blue},
emph={access,and,break,class,continue,def,del,elif ,else,
except,exec,finally,for,from,global,if,import,in,i s,
lambda,not,or,pass,print,raise,return,try,while},
emphstyle=\color{black}\bfseries,
emph={[2]True, False, None, self},
emphstyle=[2]\color{black},
emph={[3]from, import, as},
emphstyle=[3]\color{blue},
upquote=true,
columns=flexible, % pour empecher d'avoir un espacement mono
morecomment=[s]{"""}{"""},
commentstyle=\color{Hellbraun}\slshape, 
%emph={[4]1, 2, 3, 4, 5, 6, 7, 8, 9, 0},
emphstyle=[4]\color{blue},
literate=*{:}{{\textcolor{blue}:}}{1}
{=}{{\textcolor{blue}=}}{1}
{-}{{\textcolor{blue}-}}{1}
{+}{{\textcolor{blue}+}}{1}
{*}{{\textcolor{blue}*}}{1}
{!}{{\textcolor{blue}!}}{1}
{(}{{\textcolor{blue}(}}{1}
{)}{{\textcolor{blue})}}{1}
{[}{{\textcolor{blue}[}}{1}
{]}{{\textcolor{blue}]}}{1}
{<}{{\textcolor{blue}<}}{1}
{>}{{\textcolor{blue}>}}{1}
{COMPLETER}{{\textcolor{red}COMPLETER}}{1},
literate=%
            {é}{{\'{e}}}1
            {è}{{\`{e}}}1
            {ê}{{\^{e}}}1
            {ë}{{\¨{e}}}1
            {û}{{\^{u}}}1
            {ù}{{\`{u}}}1
            {â}{{\^{a}}}1
            {à}{{\`{a}}}1
            {î}{{\^{i}}}1
            {ç}{{\c{c}}}1
            {Ç}{{\c{C}}}1
            {É}{{\'{E}}}1
            {Ê}{{\^{E}}}1
            {À}{{\`{A}}}1
            {Â}{{\^{A}}}1
            {Î}{{\^{I}}}1, % pour regler les pb d accent utf8 dans les codes
%framexleftmargin=1mm, framextopmargin=1mm, frame=shadowbox, rulesepcolor=\color{blue},#1
%backgroundcolor=\color{SourceHintergrund}, 
%framexleftmargin=1mm, framexrightmargin=1mm, framextopmargin=1mm, frame=single, framerule=1pt, rulecolor=\color{black},#1
}}{}


% Définition des booleéns
\newif\iffiche
\newif\ifprof
\newif\iftd
\newif\ifcours

%%%%%%%%%%%%
% Définition des vecteurs 
%%%%%%%%%%%%
 \newcommand{\vect}[1]{\overrightarrow{#1}}
\newcommand{\axe}[2]{\left(#1,\vect{#2}\right)}

\newcommand{\rep}[1]{\mathcal{R}_{#1}}
\newcommand{\vx}[1]{\vect{x_{#1}}}
\newcommand{\vy}[1]{\vect{y_{#1}}}
\newcommand{\vz}[1]{\vect{z_{#1}}}

%%%%%%%%%%%%
% Définition des torseurs 
%%%%%%%%%%%%

 \newcommand{\torseur}[1]{%
\left\{{#1}\right\}
}

\newcommand{\torseurcin}[3]{%
\left\{\mathcal{#1} \left(#2/#3 \right) \right\}
}

\newcommand{\torseurstat}[3]{%
\left\{\mathcal{#1} \left(#2\rightarrow #3 \right) \right\}
}

 \newcommand{\torseurc}[8]{%
%\left\{#1 \right\}=
\left\{
{#1}
\right\}
 = 
\left\{%
\begin{array}{cc}%
{#2} & {#5}\\%
{#3} & {#6}\\%
{#4} & {#7}\\%
\end{array}%
\right\}_{#8}%
}

 \newcommand{\torseurcol}[7]{
\left\{%
\begin{array}{cc}%
{#1} & {#4}\\%
{#2} & {#5}\\%
{#3} & {#6}\\%
\end{array}%
\right\}_{#7}%
}

 \newcommand{\torseurl}[3]{%
%\left\{\mathcal{#1}\right\}_{#2}=%
\left\{%
\begin{array}{l}%
{#1} \\%
{#2} %
\end{array}%
\right\}_{#3}%
}

 \newcommand{\vectv}[3]{%
\vect{V\left( {#1} \in {#2}/{#3}\right)}
}


\newcommand{\vectf}[2]{%
\vect{R\left( {#1} \rightarrow {#2}\right)}
}

\newcommand{\vectm}[3]{%
\vect{\mathcal{M}\left( {#1}, {#2} \rightarrow {#3}\right)}
}


 \newcommand{\vectg}[3]{%
\vect{\Gamma \left( {#1} \in {#2}/{#3}\right)}
}

 \newcommand{\vecto}[2]{%
\vect{\Omega\left( {#1}/{#2}\right)}
}
% }$$\left\{\mathcal{#1} \right\}_{#2} =%
% \left\{%
% \begin{array}{c}%
%  #3 \\%
%  #4 %
% \end{array}%
% \right\}_{#5}}

\newcommand{\llbr}{\ensuremath{\llbracket}}
\newcommand{\rrbr}{\ensuremath{\rrbracket}}

\fichetrue
%\fichefalse

%\proftrue
\proffalse

\tdtrue
%\tdfalse

%\courstrue
\coursfalse

\newcommand{\bfsf}[1]{\textbf{\textsl{#1}}}

% -------------------------------------
% Déclaration des titres
% -------------------------------------

\def\discipline{Informatique}
\def\xxtete{Informatique}
\def\classe{PSI$\star$}
\def\xxnumpartie{Partie 5}
\def\xxpartie{Algorithmique \& Programmation II}

\def\xxnumchapitre{Chapitre 3}
\def\xxchapitre{\hspace{.12cm} Tris}

\def\xxposongletx{2}
\def\xxposonglettext{1.45}
\def\xxposonglety{13}%10

\def\xxonglet{Part. 5 -- Ch. 3}

\def\xxactivite{TD 3}
\def\xxauteur{\textsl{Serge Bays} \\ \textsl{Anthony Meurdefroid}}

\def\xxcompetences{%
\textsl{%
\textbf{Savoirs et compétences :}
\begin{itemize}[label=\ding{112},font=\color{ocre}] 
\item Alg -- C17 : tris d’un tableau à une dimension de valeurs numériques (tri par insertion, tri rapide, tri fusion).
\end{itemize}
}}

\def\xxfigures{}%figues de la page de garde

\def\xxpied{%
Partie 5 -- Algorithmique et Programmation\\
Ch 3 : Tris -- \xxactivite%
}

\def\xxtitreexo{Travaux dirigés}
\def\xxsourceexo{Xavier Pessoles \&  Skander Zannad}%\hspace{.2cm} Informatique pour tous en CPGE -- \textit{Wack \& 
% al.}}

\def\espacebandeautitre{0cm}
\setcounter{secnumdepth}{5}
%---------------------------------------------------------------------------
\begin{document}
\pagestyle{empty}


%%%%%%%% PAGE DE GARDE COURS
\ifcours
\begin{tikzpicture}[remember picture,overlay]
\node at (current page.north west)
{\begin{tikzpicture}[remember picture,overlay]
\node[anchor=north west,inner sep=0pt] at (0,0) {\includegraphics[width=\paperwidth]{\thechapterimage}};
\draw[anchor=west] (-2cm,-8cm) node [line width=2pt,rounded corners=15pt,draw=ocre,fill=white,fill opacity=0.6,inner sep=40pt]{\strut\makebox[22cm]{}};
\draw[anchor=west] (1cm,-8cm) node {\huge\sffamily\bfseries\color{black} %
\begin{minipage}{1cm}
\rotatebox{90}{\LARGE\sffamily\textsc{\color{ocre}\textbf{\xxnumpartie}}}
\end{minipage} \hfill
\begin{minipage}[c]{14cm}
\begin{titrepartie}
\begin{flushright}
\renewcommand{\baselinestretch}{1.1} 
\Large\sffamily\textsc{\textbf{\xxpartie}}
\renewcommand{\baselinestretch}{1} 
\end{flushright}
\end{titrepartie}
\end{minipage} \hfill
\begin{minipage}[c]{3.5cm}
{\large\sffamily\textsc{\textbf{\color{ocre} \discipline}}}
\end{minipage} 
 };
\end{tikzpicture}};
\end{tikzpicture}


\begin{tikzpicture}[overlay]
\node[shape=rectangle, 
      rounded corners = .25 cm,
	  draw= ocre,
	  line width=2pt, 
	  fill = ocre!10,
	  minimum width  = 2.cm,
	  minimum height = 2.5cm,] at (18cm,-4.7cm) {};
\node at (17.7cm,-4.65) {\rotatebox{90}{\textbf{\Large\color{ocre}{\classe}}}};
%{};
\end{tikzpicture}

\vspace{3.5cm}

\begin{tikzpicture}[remember picture,overlay]
\draw[anchor=west] (-2cm,-6cm) node {\huge\sffamily\bfseries\color{black} %
\begin{minipage}{2cm}
\begin{center}
\LARGE\sffamily\textsc{\color{ocre}\textbf{\xxactivite}}
\end{center}
\end{minipage} \hfill
\begin{minipage}[c]{15cm}
\begin{titrechapitre}
\renewcommand{\baselinestretch}{1.1} 
\Large\sffamily\textsc{\textbf{\xxnumchapitre}}

\Large\sffamily\textsc{\textbf{\xxchapitre}}
\vspace{.5cm}

\renewcommand{\baselinestretch}{1} 
\normalsize\normalfont
\xxcompetences
\end{titrechapitre}
\end{minipage}  };
\end{tikzpicture}
\vfill

\begin{flushright}
\begin{minipage}[c]{.3\linewidth}
\begin{center}
\xxfigures
\end{center}
\end{minipage}\hfill
\begin{minipage}[c]{.6\linewidth}
\startcontents
\printcontents{}{1}{}
\end{minipage}
\end{flushright}

\begin{tikzpicture}[remember picture,overlay]
\draw[anchor=west] (4.5cm,-.7cm) node {
\begin{minipage}[c]{.2\linewidth}
\begin{flushright}
\includegraphics[width=2cm]{png/logoCC}
\end{flushright}
\end{minipage}
\begin{minipage}[c]{.2\linewidth}
\textsl{\xxauteur} \\
\textsl{\classe}
\end{minipage}
 };
\end{tikzpicture}
\newpage
\pagestyle{fancy}

\newpage
\pagestyle{fancy}

\else
\fi


%%%%%%%% PAGE DE GARDE TD
\iftd
%\begin{tikzpicture}[remember picture,overlay]
%\node at (current page.north west)
%{\begin{tikzpicture}[remember picture,overlay]
%\draw[anchor=west] (-2cm,-3.25cm) node [line width=2pt,rounded corners=15pt,draw=ocre,fill=white,fill opacity=0.6,inner sep=40pt]{\strut\makebox[22cm]{}};
%\draw[anchor=west] (1cm,-3.25cm) node {\huge\sffamily\bfseries\color{black} %
%\begin{minipage}{1cm}
%\rotatebox{90}{\LARGE\sffamily\textsc{\color{ocre}\textbf{\xxnumpartie}}}
%\end{minipage} \hfill
%\begin{minipage}[c]{13.5cm}
%\begin{titrepartie}
%\begin{flushright}
%\renewcommand{\baselinestretch}{1.1} 
%\Large\sffamily\textsc{\textbf{\xxpartie}}
%\renewcommand{\baselinestretch}{1} 
%\end{flushright}
%\end{titrepartie}
%\end{minipage} \hfill
%\begin{minipage}[c]{3.5cm}
%{\large\sffamily\textsc{\textbf{\color{ocre} \discipline}}}
%\end{minipage} 
% };
%\end{tikzpicture}};
%\end{tikzpicture}

%%%%%%%%%% PAGE DE GARDE TD %%%%%%%%%%%%%%%
%\begin{tikzpicture}[overlay]
%\node[shape=rectangle, 
%      rounded corners = .25 cm,
%	  draw= ocre,
%	  line width=2pt, 
%	  fill = ocre!10,
%	  minimum width  = 2.5cm,
%	  minimum height = 2.5cm,] at (18.5cm,0) {};
%\node at (17.7cm,0) {\rotatebox{90}{\textbf{\Large\color{ocre}{\classe}}}};
%%{};
%\end{tikzpicture}

% PARTIE ET CHAPITRE
%\begin{tikzpicture}[remember picture,overlay]
%\draw[anchor=west] (-1cm,-2.1cm) node {\large\sffamily\bfseries\color{black} %
%\begin{minipage}[c]{15cm}
%\begin{flushleft}
%\xxnumchapitre \\
%\xxchapitre
%\end{flushleft}
%\end{minipage}  };
%\end{tikzpicture}

% Bandeau titre exo
\vspace*{\espacebandeautitre}
\begin{tikzpicture}[remember picture,overlay]
\draw[anchor=west] (-2cm,-6cm) node {\huge\sffamily\bfseries\color{black} %
\begin{minipage}{5cm}
\begin{center}
\LARGE\sffamily\color{ocre}\textbf{\textsc{\xxactivite}}

\begin{center}
\xxfigures
\end{center}

\end{center}
\end{minipage} \hfill
\begin{minipage}[c]{12cm}
\begin{titrechapitre}
\renewcommand{\baselinestretch}{1.1} 
\large\sffamily\textbf{\textsc{\xxtitreexo}}

\small\sffamily{\textbf{\textit{\color{black!70}\xxsourceexo}}}
\vspace{.5cm}

\renewcommand{\baselinestretch}{1} 
\normalsize\normalfont
\xxcompetences
\end{titrechapitre}
\end{minipage}  };
\end{tikzpicture}

\else
\fi


%%%%%%%% PAGE DE GARDE FICHE
\iffiche
\begin{tikzpicture}[remember picture,overlay]
\node at (current page.north west)
{\begin{tikzpicture}[remember picture,overlay]
\draw[anchor=west] (-2cm,-3.25cm) node [line width=2pt,rounded corners=15pt,draw=ocre,fill=white,fill opacity=0.6,inner sep=40pt]{\strut\makebox[22cm]{}};
\draw[anchor=west] (1cm,-3.25cm) node {\huge\sffamily\bfseries\color{black} %
\begin{minipage}{1cm}
\rotatebox{90}{\LARGE\sffamily\textsc{\color{ocre}\textbf{\xxnumpartie}}}
\end{minipage} \hfill
\begin{minipage}[c]{14cm}
\begin{titrepartie}
\begin{flushright}
\renewcommand{\baselinestretch}{1.1} 
\large\sffamily\textsc{\textbf{\xxpartie} \\} 

\vspace{.2cm}

\normalsize\sffamily\textsc{\textbf{\xxnumchapitre -- \xxchapitre}}
\renewcommand{\baselinestretch}{1} 
\end{flushright}
\end{titrepartie}
\end{minipage} \hfill
\begin{minipage}[c]{3.5cm}
{\large\sffamily\textsc{\textbf{\color{ocre} \discipline}}}
\end{minipage} 
 };
\end{tikzpicture}};
\end{tikzpicture}


\begin{tikzpicture}[overlay]
\node[shape=rectangle, 
      rounded corners = .25 cm,
	  draw= ocre,
	  line width=2pt, 
	  fill = ocre!10,
	  minimum width  = 2.5cm,
	  minimum height = 2.5cm,] at (18.5cm,-.6cm) {};
\node at (17.8cm,-.6cm) {\rotatebox{90}{\textsf{\textbf{\large\color{ocre}{\classe}}}}};
%{};
\end{tikzpicture}



\else
\fi



\vspace{5cm}
\pagestyle{fancy}
\thispagestyle{plain}

\setcounter{secnumdepth}{5}
\def\columnseprulecolor{\color{ocre}}
\setlength{\columnseprule}{0.4pt} 
%
%\begin{center}
%\begin{tabular}{|p{\linewidth}|}
%\hline
%\begin{center}
%\textbf{Pour l'ensemble de ce TD, il sera nécessaire de télécharger et décompresser le fichier 
%\url{https://goo.gl/cxjnS3} (aussi disponible via le site de la classe).}
%\end{center}\\
%\hline
%\end{tabular}
%\end{center}

\ifprof
\else
\begin{multicols}{2}
\fi
%---------------------------------------------------------------------------


\section*{Exercice 1 -- Ordre lexicographique}
\textit{http://mathematice.fr/ -- Serge Bays}

L’objectif est d’écrire un programme qui trie une liste de mots et les range dans l’ordre lexicographique (ordre des dictionnaires).
\begin{enumerate}
\item Écrire la définition de la variable "alphabet" :
\texttt{alphabet="AaàBbCcDdEeéèFfGgHhIiJjKkLl} \texttt{MmNnOoPpQqRrSsTtUuùVvWwXxYyZz"}.
\item Écrire une fonction \texttt{ordre\_alphabetique(c1,c2)} qui prend en argument deux caractères alphabétiques et renvoie -1 si \texttt{c1} est avant \texttt{c2}, 1 si \texttt{c2} est avant \texttt{c1} et 0 si \texttt{c1 = c2}. On pourra utiliser la
méthode index qui renvoie l’indice d’un élément dans une chaîne de caractère.
\item Écrire une fonction 

\texttt{ordre\_lexicographique(m1,m2)} qui prend en argument deux mots et renvoie -1 si "\texttt{m1 < m2}" pour l’ordre lexicographique, 0 si "\texttt{m1 = m2}" et 1 si "\texttt{m1 > m2}". On utilisera la
fonction \texttt{ordre\_alphabetique}.
\item Écrire une fonction \texttt{tri\_lexicographique(liste)} qui prend an argument une liste de mots et renvoie la liste triée. On utilisera la fonction \texttt{ordre\_alphabetique} et l’algorithme du tri par insertion.
\end{enumerate}

\ifprof
\begin{corrige}~\\

\begin{python}
alphabet='AaàBbCcDdEeéèFfGgHhIiJjKkLlMmNnOoPpQqRrSsTtUuùVvWwXxYyZz'
def ordre_alphabetique(c1,c2):
    if alphabet.index(c1)<alphabet.index(c2):
        return -1
    elif alphabet.index(c2)<alphabet.index(c1):
        return 1
    else:
        return 0

def ordre_lexicographique(m1,m2):
    n=min(len(m1),len(m2))
    for i in range(n):
        if ordre_alphabetique(m1[i],m2[i])==-1:
            return -1
        elif ordre_alphabetique(m1[i],m2[i])==1:
            return 1
    if len(m1)<len(m2):
        return -1
    elif len(m2)<len(m1):
        return 1
    else:
        return 0

def tri_lexicographique(liste):
    for i in range(len(liste)-1):
        k = i+1 # indice de la cle
        cle=liste[k]
        while ordre_lexicographique(cle,liste[k-1])==-1 and k>0:
            liste[k] = liste[k-1]
            k = k-1
        liste[k]=cle
        print("etape",i," :\t",liste)
    return liste

t=['moi', 'toi', 'bonjour','salut','bon']
print("liste :\t\t",t)
print("resultat : \t",tri_lexicographique(t))
\end{python}
\end{corrige}
\else
\fi


\section*{Exercice 2 : Problème de maçon}
\textit{D'après ressources d'Anthony Meurdefroid.}


On cherche à construire un mur à partir de briques horizontales de taille $2\times 1$ et $3\times 1$. Un mur est
correctement construit si la jointure verticale entre deux briques ne se trouve jamais immédiatement
au dessus d'une autre jointure verticale. Ainsi le mur suivant (de hauteur $h=3$ et de longueur $l=11$)%, figure~\ref{fig:mur1}) 
est correctement construit. En revanche, le mur d'après ne l'est pas.


%%%%%%%%%%%
% dessin des briques
\def\unite{0.6cm}
\tikzset{brique2/.style={anchor=south west,rectangle, rounded corners =1pt,minimum height = \unite,fill=red!10,minimum width={2*\unite}}}
\tikzset{brique3/.style={anchor=south west,rectangle, rounded corners =1pt,minimum height = \unite,fill=red!30,minimum width={3*\unite}}}
%%%%%%%%%%%%%


%\begin{figure}[h]
%\centering
%\begin{subfigure}[t]{0.45\textwidth}
%\centering
\begin{center}
\begin{tikzpicture}
\node[brique2,draw] at (0,0) {};
\node[brique3,draw] at ({2*\unite},0) {};
\node[brique2,draw] at ({5*\unite},0) {};
\node[brique2,draw] at ({7*\unite},0) {};
\node[brique2,draw] at ({9*\unite},0) {};
% deuxième couche
\node[brique3,draw] at (0,\unite) {};
\node[brique3,draw] at ({3*\unite},\unite) {};
\node[brique2,draw] at ({6*\unite},\unite) {};
\node[brique3,draw] at ({8*\unite},\unite) {};
% troisième couche
\node[brique2,draw] at (0,{2*\unite}) {};
\node[brique2,draw] at ({2*\unite},{2*\unite}) {};
\node[brique3,draw] at ({4*\unite},{2*\unite}) {};
\node[brique2,draw] at ({7*\unite},{2*\unite}) {};
\node[brique2,draw] at ({9*\unite},{2*\unite}) {};
\end{tikzpicture}
%\caption{Exemple de mur valide}

\textit{Exemple de mur valide}
\end{center}
%\label{fig:mur1}
%\end{subfigure}
%~
%\begin{subfigure}[t]{0.45\textwidth}
%\centering

\begin{center}
\begin{tikzpicture}
\node[brique2,draw] at (0,0) {};
\node[brique3,draw] at ({2*\unite},0) {};
\node[brique2,draw] at ({5*\unite},0) {};
\node[brique2,draw] at ({7*\unite},0) {};
\node[brique2,draw] at ({9*\unite},0) {};
% deuxième couche
\node[brique3,draw] at (0,\unite) {};
\node[brique3,draw] at ({3*\unite},\unite) {};
\node[brique2,draw] at ({6*\unite},\unite) {};
\node[brique3,draw] at ({8*\unite},\unite) {};
% troisième couche
\node[brique2,draw] at (0,{2*\unite}) {};
\node[brique2,draw] at ({2*\unite},{2*\unite}) {};
\node[brique2,draw] at ({4*\unite},{2*\unite}) {};
\node[brique3,draw] at ({6*\unite},{2*\unite}) {};
\node[brique2,draw] at ({9*\unite},{2*\unite}) {};
% mise en évidence
\draw[ultra thick, blue] ({6*\unite},{3*\unite}) -- ({6*\unite},{1*\unite});

\end{tikzpicture}

%\caption{Exemple de mur non valide}
\textit{Exemple de mur non valide}
\end{center}
%\label{fig:mur2}
%\end{subfigure}
%\caption{Exemples de mur de hauteur $h=3$ et de longueur $l=11$}
%\label{fig:murs}
%\end{figure}

L'objectif de ce problème est de dénombrer le nombre de façons différentes de construire un mur de
longueur $l$ et de hauteur $h$.




%%%%%%%%%%%%%%%%%%%%%%%%%%%%%%%%%

%% Codage et fonctions

%%%%%%%%%%%%%%%%%%%%%%%%%%%%%%%%%

\subsection*{Codage du problème et fonctions utiles}
\label{sec:codageFonctions}

Chaque rangée de briques (une ligne) de longueur $l$ va être codée comme une chaîne de caractères, composée uniquement de \texttt{'0'} et \texttt{'1'}, de longueur $l-1$. Par exemple, la rangée de taille 11 %(Fig.~\ref{fig:rangeeDeBriques}) 
est codée dans une chaîne de caractères de longueur 10.%
% (Fig.~\ref{fig:exempleCodage}).




%\begin{figure}[h]
%\centering
%\begin{subfigure}[t]{0.45\textwidth}
%\centering

\begin{center}

\begin{tikzpicture}

\node[brique2,draw] at (0,0) {};
\node[brique3,draw] at ({2*\unite},0) {};
\node[brique2,draw] at ({5*\unite},0) {};
\node[brique2,draw] at ({7*\unite},0) {};
\node[brique2,draw] at ({9*\unite},0) {};
\end{tikzpicture}

%\caption{Exemple de rangée de briques}
\textit{{Exemple de rangée de briques}}
\end{center}
%\label{fig:rangeeDeBriques}
%\end{subfigure}
%~
%\begin{subfigure}[t]{0.45\textwidth}
%\centering
\begin{center}
\begin{tikzpicture}
\def\uniteCodage{0.4cm}
\tikzset{codage/.style={anchor=south west,rectangle,minimum width=\uniteCodage},minimum height = \uniteCodage};
\node[codage,draw] at (0,0) {\texttt{'0100101010'}};
\end{tikzpicture}

\textit{{Codage de la rangée de briques}}
\end{center}
%\caption{Codage de la rangée de briques}
%\label{fig:exempleCodage}
%\end{subfigure}
%\caption{Illustration du codage choisi}
%\label{fig:codage}
%\end{figure}

Dans le détail, le \textit{i}\up{ième} élément (chiffre) de la chaîne de caractère indique si au niveau de la \textit{i}\up{ième} unité de longueur il y a oui ou non une jointure entre deux briques (\texttt{'1'} si il y a une jointure, \texttt{'0'} sinon). On n'inclut pas les deux extrémités - ce qui donne :

%\begin{figure}[h]
%\centering
\begin{center}
\begin{tikzpicture}
\node[brique2,draw] at (0,0) {};
\node[brique3,draw] at ({2*\unite},0) {};
\node[brique2,draw] at ({5*\unite},0) {};
\node[brique2,draw] at ({7*\unite},0) {};
\node[brique2,draw] at ({9*\unite},0) {};

% valeurs
\def\uniteCodage{0.4cm}
\tikzset{valeur/.style={anchor=north}};
\node[valeur] (N1) at ({1*\unite},{-1*\uniteCodage}) {\texttt{'0}};
\node[valeur] (N2) at ({2*\unite},{-1*\uniteCodage}) {\texttt{1}};
\node[valeur] (N3) at ({3*\unite},{-1*\uniteCodage}) {\texttt{0}};
\node[valeur] (N4) at ({4*\unite},{-1*\uniteCodage}) {\texttt{0}};
\node[valeur] (N5) at ({5*\unite},{-1*\uniteCodage}) {\texttt{1}};
\node[valeur] (N6) at ({6*\unite},{-1*\uniteCodage}) {\texttt{0}};
\node[valeur] (N7) at ({7*\unite},{-1*\uniteCodage}) {\texttt{1}};
\node[valeur] (N8) at ({8*\unite},{-1*\uniteCodage}) {\texttt{0}};
\node[valeur] (N9) at ({9*\unite},{-1*\uniteCodage}) {\texttt{1}};
\node[valeur] (N10) at ({10*\unite},{-1*\uniteCodage}) {\texttt{0'}};

%relier
\draw[>=stealth,->] (N1.north) -- ({1*\unite},0);
\draw[>=stealth,->] (N2.north) -- ({2*\unite},0);
\draw[>=stealth,->] (N3.north) -- ({3*\unite},0);
\draw[>=stealth,->] (N4.north) -- ({4*\unite},0);
\draw[>=stealth,->] (N5.north) -- ({5*\unite},0);
\draw[>=stealth,->] (N6.north) -- ({6*\unite},0);
\draw[>=stealth,->] (N7.north) -- ({7*\unite},0);
\draw[>=stealth,->] (N8.north) -- ({8*\unite},0);
\draw[>=stealth,->] (N9.north) -- ({9*\unite},0);
\draw[>=stealth,->] (N10.north) -- ({10*\unite},0);

\end{tikzpicture}
\end{center}
%\end{figure}


\subparagraph{}\textit{Dessiner la rangée de briques correspondant à la chaîne de caractère suivante : \texttt{'0100100'}. De quelle taille est cette rangée ?}

%%\reponsebox{1.75cm}

%%\corrigebox{
\ifprof
\begin{corrige}~\\
\begin{center}
\begin{tikzpicture}
\def\unite{0.6cm}
\tikzset{brique/.style={anchor=south west,rectangle, rounded corners =1pt,minimum height = \unite,fill=red!10}};
\node[brique,minimum width={2*\unite},draw] at (0,0) {};
\node[brique,minimum width={3*\unite},draw] at ({2*\unite},0) {};
\node[brique,minimum width={3*\unite},draw] at ({5*\unite},0) {};
\end{tikzpicture}
\end{center}

La rangée de briques est : 1 brique de taille 2, puis 2 de taille 3, soit une rangée de
taille \fbox{8}.%}
\end{corrige}

\else
\fi

%Maintenant, on suppose fixée à l'avance la longueur des murs que l'on souhaite compter. On appellera cette variable \texttt{tailleMax}.


\subparagraph{}\textit{\'Ecrire une fonction \texttt{imprimeRangee(rangee)} qui prend en argument une liste des deux chiffres 2 et 3 correspondant à une rangée de briques et qui renvoie la rangée de briques sous la forme codée (type \texttt{string}). Exemple : \texttt{imprimeRangee([2,3,2,2,2])} renvoie \texttt{'0100101010'}.}

%\reponsebox{4cm}
\ifprof
\begin{corrige}
\begin{python}
def imprimeRangee(rangee):
    '''rangee est une liste de 2 et de 3'''
    c=''
    for element in rangee:
        if element==2:
            c+='01'
        else:
            c+='001'
    return c[:-1]
\end{python}
\end{corrige}
\else
\fi
%\corrigebox{\pythonexternal{imprimeRangee.py}}


Deux rangées sont dites compatibles si elles peuvent être placées l'une sur l'autre, c'est-à-dire que les jointures ne sont pas au même endroit.

\subparagraph{}\textit{\'Ecrire une fonction}

\textit{ \texttt{sontCompatibles(r1,r2)} prenant en argument \texttt{r1} et \texttt{r2}, 2 rangées de briques sous la forme codée (type \texttt{string} - ex : \texttt{'011100'}), et qui renvoie \texttt{True} si les deux rangées données sont compatibles, \texttt{False} sinon.}

%\reponsebox{5cm}

%\corrigebox{\pythonexternal{sontCompatibles.py}}
\ifprof
\begin{corrige}~\\

\begin{python}
def sontCompatibles(r1,r2):
    for i in range(len(r1)):
        if r1[i]=='1' and r2[i]=='1':
            return(False)
    return(True)
\end{python}
\end{corrige}
\else
\fi

%%%%%%%%%%%%%%%%%%%%%%%%%%%%%%%%%%%%%%%%%
%%%%%%%%%%% resolution %%%%%%%%%%%%%%%%%%
%%%%%%%%%%%%%%%%%%%%%%%%%%%%%%%%%%%%%%%%%

\subsection*{Résolution du problème}
\label{sec:resolution}



Pour dénombrer le nombre de façons différentes de construire un mur de longueur $l$ et de hauteur $h$, on adopte le principe de résolution suivant :
\begin{itemize}
\item \textbf{phase 1 :} déterminer toutes les combinaisons possibles des briques permettant de construire une rangée de longueur $l$. Par exemple, il n'y a que deux façons possibles pour les rangées de longueur 5 : une brique de longueur 2, puis une brique de longueur 3 ; et l'inverse. Les codages correspondants sont : \texttt{'0100'} et \texttt{'0010'}. Soit $R$ l'ensemble de ces codages de rangées de briques;
\item \textbf{phase 2 :} pour chaque rangée $r$ de $R$, déterminer le nombre de murs de hauteurs $h$ dont la première rangée de briques (celle tout en bas) est $r$. Cette valeur est notée \texttt{compteRangee(R, r, h)};
\item \textbf{phase 3 :} calculer $\sum \limits_{r \in R} \mathtt{compteRangee(R,r,h)}$ pour toutes les rangées possibles de longueur $l$. C'est le résultat final attendu.
\end{itemize}


%Dans la suite, l'objectif est de comprendre l'algorithme récursif de la phase 2, puis de l'améliorer.


\subsection*{\'Etude de la phase 1}
\label{ssec:phase1}

Dans un premier temps, on aimerait connaître le nombre de combinaisons possibles distinctes afin de remplir la rangée de briques avec des briques élémentaires de dimensions 2 et 3.

\subparagraph{}\textit{Proposer une fonction récursive \texttt{combinaisons(l)} prenant en argument la longueur du mur $l$ et renvoyant le nombre de solutions possibles distinctes afin de remplir une rangée. Vous pourrez construire votre algorithme en vous appuyant sur un schéma (arbre) ayant pour sommet $l=11$.}
%\reponsebox{9cm}

%\corrigebox{\pythonexternal{combinaisons.py}}
\ifprof
\begin{corrige}~\\

\begin{python}
def combinaisons(n):
    if n==1:
        return(0)
    elif n==2 or n==3:
        return(1)
    else:
        return(combinaisons(n-2)+combinaisons(n-3))
\end{python}
\end{corrige}
\else
\fi



\subparagraph{}\textit{Préciser la complexité de votre fonction.}
%\reponsebox{4 cm}
%\corrigebox{
\ifprof
\begin{corrige}
On a : $C(n)=C(n-2)+C(n-3)$, alors on obtient une complexité en $C(n)=O\left(\lambda^n \times 2^{n}\right)$ soit \fbox{une complexité exponentielle}.

\end{corrige}

Pour la phase 1, on suppose écrite une fonction \texttt{remplissage(l)} qui renvoie la liste de toutes les rangées possibles de longueur $l$. Par exemple, \texttt{remplissage(5)} renvoie les deux rangées : \texttt{'0100'} et \texttt{'0010'}.








%Le code de la fonction toutes les lignes vous est donné blabla (à décider..)

\subparagraph{}\textit{Sans écrire aucun code, quelle serait la structure du code de la fonction \texttt{remplissage} ?}

%\reponsebox{6cm}

%\corrigebox{
\ifprof
\begin{corrige}~\\
Il est nécessaire de créer une fonction récursive qui ajoute dans la branche de gauche une brique de 2 (concaténation de la chaîne de caractère en cours avec \texttt{'10'}), puis dans la branche de droite une brique de trois (concaténation de la chaîne de caractère en cours avec \texttt{'100'}). Il faut gérer la longueur du mur en cours (\texttt{w} dans le code ci-dessous) - ceci constituera notre cas de base dès lors que le mur en cours fera la longueur souhaitée. Enfin, il faut remplir une liste de toutes les combinaisons. J'ai utilisé une variable globale (très utile ici), en s'appuyant sur la génération pour une brique initialement de 2 puis une de 3. Le code à titre informatif:

\begin{python}
def remplissage(w,r,tailleMax):
    if w==tailleMax:
        toutesLesCombis.append(r)
    elif w<tailleMax:
        remplissage(w+2,r+'10',tailleMax)
        remplissage(w+3,r+'100',tailleMax)
        
global toutesLesCombis
toutesLesCombis=[]

tailleMax=9
      
remplissage(2,'0',tailleMax)
remplissage(3,'00',tailleMax)
\end{python}
%\pythonexternal{remplissage.py}
 %}
 \end{corrige}
\else
\fi

\subsection*{\'Etude de la phase 2}
\label{ssec:phase2}

Pour dénombrer les murs (phase 2 de la résolution), on utilise une fonction récursive \texttt{compteRangee(R, r, h)}
qui détermine le nombre de murs dont la première rangée de briques est $r$ et dont la hauteur est $h$.
L'objet de cette sous-section est de comprendre cet algorithme et son implémentation, qui sont donnés dans le fichier macon.py%en page~\pageref{LastPage}.

\subparagraph{}\textit{Pourquoi la quantité \texttt{compteRangee(R, r, 1)} vaut-elle 1 pour tout $r$ ?}
%\reponsebox{3cm}

%\corrigebox{}

\ifprof
\begin{corrige}
Parce qu'il n'y a qu'une seule façon de faire un mur de hauteur 1 avec une rangée donnée !
\end{corrige}
\else
\fi


\subparagraph{}\textit{On suppose $l = 5$. Développer sur votre copie les différentes étapes de calcul (avec les
appels récursifs) de l'appel \texttt{compteRangee(combinaisonsCinq,combinaisonsCinq[0],2)}, \texttt{combinaisonsCinq} étant la liste comportant les deux chaînes de caractères, \texttt{'0100'} et \texttt{'0010'}, codant les rangées.
}
%\reponsebox{8cm}

%\corrigebox{
\ifprof
\begin{corrige}
On va donc dérouler les appels récursifs (arbre) :
\begin{itemize}
\item appel à \texttt{compteRangee(combinaisonsCinq,'0100',2)}
\item initialement, \texttt{somme} vaut 0 ; nbCombi vaut 2 ; i vaut 0
\item i<2 ici 0<2 vrai
\item SontCompatibles('0100','0100') False ; i=i+1=1
\item i<2 ici 1<2 vrai
\item SontCompatibles('0010','0100') True
\begin{itemize}
\item appel à \texttt{compteRangee(combinaisonsCinq,'0010',1)}
\item h=1 ; la fonction renvoie 1 ; somme=somme+1=1
\item i=2 ; boucle while interrompue
\end{itemize}
\item à la fin \texttt{somme} vaut 1, et on renvoie 1.
\end{itemize}
%}
\end{corrige}
\else
\fi
%\subparagraph{}\textit{Démonter en exhibant un invariant de boucle que la fonction proposée permet de renvoyer le résultat attendu. Préciser le variant. }
%
%%\reponsebox{10cm}
%
%%\corrigebox{}




\subsection*{\'Etude de la phase 3}
Afin de réaliser l'objectif, la fonction \texttt{compteTout(toutesLesCombinaisonsLigne,hauteur)} sera écrite. 


\subparagraph{}\textit{\'Ecrire la fonction \texttt{compteTout(toutesLesCombinaisonsLigne,hauteur)} prenant en argument une liste de chaînes de caractères et la hauteur du mur souhaitée, et renvoyant un entier nombre de combinaisons possibles afin de construire le mur d'une manière correcte.}
%\reponsebox{8cm}
%\corrigebox{\pythonexternal{compteTout.py}}
\ifprof
\begin{corrige}
def compteTout(toutesLesCombinaisonsLigne,hauteur):
    nbCombi=len(toutesLesCombinaisonsLigne)
    somme=0
    i=0
    while i<nbCombi:
        temp=compteRangee(toutesLesCombinaisonsLigne,toutesLesCombinaisonsLigne[i],hauteur)
        somme+=temp
        i+=1
    return(somme)
\end{corrige}
\else
\fi
Dans cette sous-section, nous chercherons les limites de notre algorithme naïf proposé dans la sous-section précédente, et nous chercherons donc à l'améliorer.

\subparagraph{}\textit{En s'appuyant sur le calcul des murs de longueur 9 et de hauteur 3, donner le principal défaut de l'algorithme naïf proposé précédemment.}

%\reponsebox{4cm}

\ifprof
\begin{corrige}
Pour le calcul des murs de longueur 9 et de hauteur 3, une bête réflexion au début de chaque appel récursif nous montre que l'on appelle deux fois \texttt{compteRangee(toutesLesCombinaisons,'00100100',2)}. Ce phénomène se répétant très souvent, il rend très lent notre algorithme.
\end{corrige}
\else
\fi


\subparagraph{}\textit{Proposer alors une amélioration de l'algorithme.}

%\reponsebox{4cm}

%\corrigebox{
\ifprof
\begin{corrige}
Si on reprend l'exemple précédent, on peut donc effectuer
ce calcul une fois, puis stocker ce nombre dans une table. Cela devient très fortement intéressant à partir du moment où l'on demande par exemple le calcul des murs de longueur 30 et de hauteur 10.
\end{corrige}
\else
\fi
%
%%
%%%\newpage
%%\begin{center}
%%\begin{Huge}
%%\fbox{Annexe}
%%\end{Huge}
%%\end{center}
%\subsection*{Annexe}
%\begin{python}
%def compteRangee(toutesLesCombinaisonsLigne,r,hauteur):
%    if hauteur==1:
%        return(1)
%    else:
%        nbCombi=len(toutesLesCombinaisonsLigne)
%        somme=0
%        i=0
%        while i<nbCombi:
%            if sontCompatibles(r,toutesLesCombinaisonsLigne[i]):
%                somme+=compteRangee(toutesLesCombinaisonsLigne,toutesLesCombinaisonsLigne[i],hauteur-1)
%            i+=1
%        return(somme)
%\end{python}



\ifprof
\else
\end{multicols}
\fi
%\begin{thebibliography}{2}
%\bibitem{1}{Patrick Beynet, \textit{Supports de cours de TSI 2}, Lycée Rouvière, Toulon.}
%\bibitem{2}{<< Mandel zool 08 satellite antenna >>. Sous licence CC BY-SA via Wikimedia Commons - 
%\url{https://fr.wikipedia.org/wiki/Ensemble_de_Mandelbrot#/media/File:Mandel_zoom_08_satellite_antenna.jpg}}
%\bibitem{3}{\url{http://www.obside.fr/fractales/pages/Recursif/}}
%\end{thebibliography}
\end{document}

